
\documentclass{resume} % Use the custom resume.cls style

\usepackage[left=0.75in,top=0.6in,right=0.75in,bottom=0.6in]{geometry} % Document margins
\newcommand{\tab}[1]{\hspace{.2667\textwidth}\rlap{#1}}
\newcommand{\itab}[1]{\hspace{0em}\rlap{#1}}
\name{Ebraheem Khaled Farag} % Your name
\address{733 Weston Park Drive, Powell Ohio} % Your address
%\address{123 Pleasant Lane \\ City, State 12345} % Your secondary addess (optional)
\address{(+1)~614~623~4875 \\ ekfarag@asu.edu} % Your phone number and email
\address{Orcid: 0000-0002-5794-4286}
\begin{document}

%----------------------------------------------------------------------------------------
%	EDUCATION SECTION
%----------------------------------------------------------------------------------------

\begin{rSection}{Education}

{\bf The Ohio State University} \hfill {\em  2014 - 2019} 
\\ B.S. \hfill { Overall GPA: 3.76}
\\ Engineering Physics, Mechanical Specialization\hfill { College of Engineering}

{\bf Arizona State University} \hfill {\em August 2019 - Present} 
\\ 5th year PhD Candidate \hfill { Overall GPA: 3.95}
\\ Astrophysics \hfill { School of Earth and Space Exploration}


\end{rSection}
%----------------------------------------------------------------------------------------
%	TECHNICAL STRENGTHS SECTION
%----------------------------------------------------------------------------------------

\begin{rSection}{Technical Strengths}

\begin{tabular}{ @{} >{\bfseries}l @{\hspace{6ex}} l }
Computer Languages &  Fortran, MATLAB, Python, C++ \\
Software & MESA, GYRE, YREC, \LaTeX, Mathematica \\
\end{tabular}

\end{rSection}

%----------------------------------------------------------------------------------------
%	WORK EXPERIENCE SECTION
%----------------------------------------------------------------------------------------

\begin{rSection}{Research Experience}

\begin{rSubsection}{Los Alamos National Laboratory}{May 2015 - present}{Graduate Student Researcher, Katie Mussack, Suzannah Wood, \& Joyce Guzik}{}
%\item Developed new solar-evolution-models to explore the solar-abundance-problem. 
\item Validating proprietary LANL atomic opacity data with modern stellar evolution models and helioseismology.
%\item Explored the effect of early mass loss, rotation, and magnetic braking on the solar structure.
\end{rSubsection}

% \begin{rSubsection}{The Ohio State University}{January 2015 - June 2015}{Student Research Assistant, Department of Physics, Chris Orban}{}
% \item Used knowledge of high energy-density systems, multiple mean-free-path calculations were computed in MATLAB; to generalize basic constraints necessary for laser shots at the National Ignition Facility (NIF) and OMEGA laser facility.

% \end{rSubsection}


\begin{rSubsection}{Arizona State University}{August 2019 - Present}{Graduate Research Assistant, School of Earth and Space Exploration (SESE), Frank Timmes}{}
\item Using the MESA Stellar evolution code to study gravitational wave detections of merging binary black holes, neutrino astronomy, massive stars, Helium burning nuclear reactions, and Pair-instability supernova.
\end{rSubsection}





\end{rSection}


% \begin{rSection}{Work Experience}

% \begin{rSubsection}{Siemens Energy}{January 2017 - June 2017}{Fluid Systems Engineer}{}
% \item Worked with the Trent 60 MegaWatt aero-derivative gas turbine.
% \item Contributed to the design of an external inlet spray intercooling mechanism for cooling intake air to the gas turbine for higher efficiency.
% \item Helped specify more than fifty high pressure valves for fuel injection into the gas turbine.
% \item Reduced gas turbine installation time, by conducting a full re-haul of installation logistics.
% \end{rSubsection}

% \end{rSection}






\begin{rSection}{Service Work}

\begin{rSubsection}{Astrophysical Journal Referee  :}{2021}{}
\item Has completed a referee report for ApJ. 

\end{rSubsection}
\end{rSection}


\begin{rSection}{Awards}

\begin{rSubsection}{AAS 237 Chambliss Honorable Mention :}{2021}{}
\item Exploring the Black Hole Mass Gap Sensitivity to $^{12}$C($\alpha$,$\gamma$)$^{16}$O. 
\end{rSubsection}
\end{rSection}




\begin{rSection}{Invited Talks}

\begin{rSubsection}{Los Alamos Distinguished Seminar Series :}{2019}{}
\item Is There A Solar Model Solution to The Faint Young Sun Paradox? 
\end{rSubsection}
\begin{rSubsection}{Los Alamos Distinguished Seminar Series :}{2022}{}
\item Resolving the Peak of the Black Hole Mass Spectrum 
\end{rSubsection}
\end{rSection}



%----------------------------------------------------------------------------------------
\begin{rSection}{Publications} \itemsep -3pt

\item -Farag, E., Renzo, M., Farmer, R., Chidester, M. T., \&
Timmes, F. X. 2022, ApJ, 937, 112,
\url{doi: 10.3847/1538-4357/ac8b83}
\item -Chidester, M. T., Farag, E., & Timmes, F. X. 2022, ApJ, 935, 21, \url{doi: 10.3847/1538-4357/ac7ec3}
\item -Chidester, M. T., Timmes, F. X., Schwab, J., et al. 2021, ApJ, 910, 24, \url{doi: 10.3847/1538-4357/abdec4}
\item Farag, E., & Timmes, F. 2021, in American Astronomical Society Meeting Abstracts, Vol. 53, American Astronomical Society Meeting Abstracts, 524.03
\item -Guzik, J. A., Farag, E., Ostrowski, J., et al. 2020, arXiv e-prints, arXiv:2002.04073.\newline https://arxiv.org/abs/2002.04073
\item -Guzik, J. A., Farag, E., Ostrowski, J., et al. 2020b, in The 35th Annual New Mexico Symposium, ed. A. D. Kapinska, 5,
\item -Guzik, J. A., Farag, E., Ostrowski, J., et al. 2020a, Journal of the American Association of Variable Star Observers (JAAVSO), 48, 102
\item -Farag, E., Timmes, F. X., Taylor, M., Patton, K. M., & Farmer, R. 2020a, ApJ, 893, 133,
\item -Farag, E. K., Timmes, F., Taylor, M., Patton, K. M., & Farmer, R. 2020b, in American Astronomical Society Meeting Abstracts, Vol. 236, American Astronomical Society Meeting Abstracts #236, 331.01
doi: 10.3847/1538-4357/ab7f2c
\item -A diffusive radiation hydrodynamics code, xRage, is implemented to compare radiation flow with experimental data from the Omega laser facility. Vandervort, Robert; Elgin, Laura; Farag, Ebraheem; Mussack, Katie; Baumgaertel, Jessica Ann; Keiter, Paul; Klein, Sallee; Orban, Christopher; Drake, R. Paul
\item -Joyce A. Guzik , C.J. Fontes , P. Walczak , S.R. Wood , K. Mussack , and E. Farag. “Sound speed and oscillation frequencies for solar models evolved with Los Alamos ATOMIC opacities”. Astronomy in Focus, Volume 2 XXIXth IAU General Assembly, August 2015 Piero Benvenuti, ed. c 2015 International Astronomical Union.













\end{rSection}

\end{document}
